\chapter{基础模型}
\label{chap:nlp-classical-models}

同一句“读书会改至周五”可以有不同的分析结果，但在计算分析之前，还有一个更基本的问题：
我们用什么数学对象描述语言规律？若给每个完整句子单独计数，几乎每个新句子都没有记录；
若把句子拆成相互独立的词，又会丢掉“改至”与新日期之间的联系。统计建模需要在保留依赖
和获得足够样本之间作出取舍。本章的“基础模型”指大语言模型普及以前形成的主要建模思路，
其中也包括早期神经语言模型，不把它用作英文foundation model的译名。

我们把建模分成四个相互配合的决定：描述哪些变量的概率，怎样分解变量间的依赖，怎样用
有限参数表达局部规律，以及怎样根据观测估计参数。这些决定不是互斥的模型类别。
同一个线性链模型既可以用概率表参数化，也可以采用神经网络计算势函数；同一个词向量又
可以服务于顺序预测或结构化判定。沿这四个决定展开，才能在后续任务变化时复用建模认识，
而不是为分类、标注和生成各背一套模型名单。

\section{语言现象的概率描述}
\label{sec:nlp-classical-probability}

\subsection{观测序列与随机对象}
\label{sec:nlp-classical-objects}

先固定词元化规则和有限词表$\mathcal V$，把文本写成$w_{1:T}$。词元可以是词、字或子词，
选择改变了样本空间，也改变了计数的含义。结束符属于序列的一部分；没有长度上界时，模型
还须以概率一结束，才能在有限字符串集合上形成归一化分布。固定长度的条件分布与任意长度
文本的分布是不同对象，不能只写一个乘积就忽略结束问题。

为了说明这些决定，以下采用人工构造的小语料：“会议／改至／周五／EOS”出现两次，
“会议／定于／周五／EOS”出现一次。我们暂时把每一行视为一次独立观测。这个假设使各行
的似然可以相乘，却没有要求一行内部的词元独立。实际重复转载的三行可能只有一个信息
来源，独立行的数学形式不能替代对语料来源的判断。

\subsection{联合分布与条件分布}
\label{sec:nlp-classical-joint-conditional}

若希望比较文本本身的可能性，可以直接描述$p(w_{1:T})$。若输入文本$x$已经给出，需要
比较的是一个标注序列$y$，则可以描述联合分布$p(x,y)$，也可以直接描述条件分布$p(y\mid x)$。
前者必须同时分配不同文本的概率；后者只要求对同一个$x$下的所有合法$y$归一化。
当$p(x)>0$时，两者通过$p(y\mid x)=p(x,y)/\sum_{y'}p(x,y')$联系，但从条件分布不能唯一
恢复输入分布。因此，二者并不是同一个学习问题的不同写法。

例如“周五”可能标作日期，也可能出现在作品名称中。联合模型要说明词与标签如何共同
出现，条件模型则可以把整段文本中的引号、位置和相邻词作为已知信息。是否使用未来词元，
取决于输入在应用时是否已经完整可见；对已收到整段通知的标注器可见，不表示实时逐词
预测也能使用。先确定可用信息，再决定模型条件，才能避免把测试时没有的线索放进计算。

\subsection{隐变量与可观测证据}
\label{sec:nlp-classical-latent}

有些依赖不能直接从词面读出。我们可以引入词性、句法树或主题等隐变量$z$，用
\begin{equation}
 p_{\bm\theta}(x)=\sum_{z\in\mathcal Z(x)}p_{\bm\theta}(x,z)
 \label{eq:nlp-classical-marginal}
\end{equation}
对所有允许的隐结构求和。$\mathcal Z(x)$是与输入相容的结构集合，$\bm\theta$是参数。
求和保留了多种解释的不确定性；只取概率最大的$z$是近似或另一个预测目标，一般不等于
求和。例如两种结构的联合概率为$0.3$和$0.2$，观测概率是$0.5$，最大项却只有$0.3$。

隐变量的名称不保证它就是人类期望的语言概念。如果交换两个隐状态的编号不改变所有观测
概率，数据就无法确定哪个状态该叫“名词”。加入标注、结构约束或先验可以减少这种
不确定性，但解释仍须由外部证据检查。这使下一步的结构选择既是计算决定，也是对可表达
规律的限制。

\section{语言依赖的结构分解}
\label{sec:nlp-classical-dependencies}

\subsection{有限历史与局部链式依赖}
\label{sec:nlp-classical-local-chain}

概率的链式法则总能把固定序列分成逐位置条件概率；计算困难在于每个条件可能包含整个
历史。$n$元模型（$n$-gram model）采用有限历史近似：
\begin{equation}
 p(w_{1:T})=\prod_{t=1}^T p(w_t\mid w_{<t})
 \;\approx\;\prod_{t=1}^T p(w_t\mid w_{\max(1,t-n+1):t-1}).
 \label{eq:nlp-classical-ngram}
\end{equation}
开头用约定的起始符补足历史，末尾预测EOS。左侧等式来自概率法则，右侧近似来自建模
假设；数据再多也不会让固定二元模型重新读取被删除的更早词元。

\term{隐马尔可夫模型}（hidden Markov model, HMM）把局部依赖放在隐状态$z_t$上。
为单独比较状态依赖，这里切换为已经分段的正文观测：$T$只计正文词元，
观测词表不含EOS，式子只描述给定长度的正文与状态。
用初始概率$\pi$、转移表$A$和发射表$B$写成
\begin{equation}
 p(w_{1:T},z_{1:T}\mid T)=\pi_{z_1}B_{z_1,w_1}
 \prod_{t=2}^T A_{z_{t-1},z_t}B_{z_t,w_t}.
 \label{eq:nlp-classical-hmm}
\end{equation}
这里显式以长度$T$为条件；若用于产生可变长完整语串，就须另建长度或结束机制。状态能汇集不相邻的
词面证据，但给定$z_t$后当前词只依赖该状态，是实质性的条件独立假设。
HMM的概率与推断基础可参考\BookChapterRef{chap:pgm-dynamic}{隐马尔可夫模型}。
本章关注这项假设如何把一个全句问题改成可反复计算的局部量。

对给定文本的标签序列，\term{线性链条件随机场}（conditional random field, CRF）保留标签链，
但允许局部势读取整段$x$：
\begin{equation}
 p_{\bm\theta}(y\mid x)=\frac{1}{Z_{\bm\theta}(x)}
 \exp\!\left\{\sum_{t=1}^{T}\bm\theta^{\mathsf T}
 \bm f(y_{t-1},y_t,x,t)\right\}.
 \label{eq:nlp-classical-crf}
\end{equation}
$\bm f$是特征向量，$y_0$是起始标签，$Z_{\bm\theta}(x)$对合法标签序列求和归一化。
它没有沿词元轴对观测作HMM式独立假设；高效链式求和依赖于标签耦合仍局限在相邻位置。
若势同时连接任意远处的多个标签，原来的动态规划就不能直接复用。

\subsection{层次组合与递归依赖}
\label{sec:nlp-classical-hierarchy}

局部链条可以描述相邻关系，却不直接表达“一个名词短语可以包含另一个修饰短语”的
嵌套结构。\term{概率上下文无关文法}（probabilistic context-free grammar, PCFG）为产生式
$A\to\beta$分配概率，并要求相同左部$A$的各规则概率之和为一。这里由文法另行规定终结符词表，
句尾由有限推导树的结束确定；若要统一到含EOS的字符串接口，可以给顶层增加
生成EOS的确定性规则。对一棵有限推导树$\tau$，其联合概率由所有规则应用相乘：
\begin{equation}
 p(w,\tau)=\prod_{r\in\tau}p(r),\qquad p(w)=\sum_{\tau:\,\operatorname{yield}(\tau)=w}p(w,\tau).
 \label{eq:nlp-classical-pcfg}
\end{equation}
树的叶序列$\operatorname{yield}(\tau)$等于观测句子。这里相同非终结符使用同一规则分布，
所以递归可以扩展到训练中未出现的树形；也正因为这种共享，规则选择看不到任意祖先与
词汇条件。递归文法还需检查终止性，局部规则归一化并不自动保证有限句子的总概率为一。

若$S\to NP\ VP$的概率为一，$NP\to$“会议”的概率为$0.6$，$VP\to$“改期”的概率
为$0.4$，则这棵特定树产生“会议改期”的概率为$0.24$。另一个合法树若产生同一句子，
句子概率必须把它加进来。应用章的句法分析将研究怎样找树；此处先明确，句子相同不代表
隐结构相同，结构概率最大与句子边缘概率最大也不是同一对象。

\subsection{交换性与混合结构}
\label{sec:nlp-classical-mixture}

有时需要保留的是文档中哪些词共同出现，而不是它们的先后顺序。\term{潜在狄利克雷分配}
（latent Dirichlet allocation, LDA）让每篇文档$d$有主题比例$\bm\vartheta_d$，每个
主题$k$有词分布$\bm\phi_k$，并令每个词元先选择主题，再选择词：
\begin{equation}
 p(w_{dn}=v\mid\bm\vartheta_d,\bm\phi)
 =\sum_{k=1}^{K}\vartheta_{dk}\phi_{kv},\qquad
 p(w_{d,1:N_d}\mid\bm\vartheta_d,\bm\phi)
 =\prod_{n=1}^{N_d}\sum_k\vartheta_{dk}\phi_{k,w_{dn}}.
 \label{eq:nlp-classical-lda}
\end{equation}
通常再为比例与主题词分布指定狄利克雷先验。给定这些量时，交换词元位置不改变概率，
这就是该模型保留的交换性。积分掉共同的文档比例后，文档内词元会相关，不能再把“给定
隐变量时独立”说成“边缘上独立”\scite{blei2003}。

交换性用失去语序换取集合统计的简化。“允许校外参加，禁止校内参加”和“允许校内参加，禁止校外参加”
在相同分词下具有同样的词袋，却把允许与禁止绑定到了不同人群；这不是调优主题数就能消除的问题。链、树和混合结构分别
保留不同依赖，任务需要哪种证据，才决定哪种简化可接受。下一节固定结构，考察局部规律
怎样用参数表达。

\section{语言规律的参数化与共享}
\label{sec:nlp-classical-parameterization}

\subsection{概率表与离散事件共享}
\label{sec:nlp-classical-tables}

二元模型可以为每个历史词$h$存一行概率$p(v\mid h)$。行内各项非负且和为一，历史完全
不同的两行则可以自由变化。词表大小为$V$时，完整二元表约有$V(V-1)$个自由参数；
$n$增大后，历史组合数量随$V^{n-1}$增长。这个数量描述候选模型的表达自由度，不意味着
所有单元都能从有限语料得到可靠估计。

把未知词映射到同一符号、把历史归入词类，或对低阶分布作插值，都会使不同事件共享信息。
例如两个罕见地点都映射为“地点”，能增加统计量，却丢失它们之间的差别。平滑是在给定
计数与约束下重新分配概率，下一节具体计算；结构聚合则直接改变了模型能区分哪些历史。
这两类操作都能缓解稀疏，但作用对象不同。

\subsection{特征组合与对数线性评分}
\label{sec:nlp-classical-loglinear}

若希望一个词的后缀、左右词和位置都参与判断，为每种组合独立建表会继续放大稀疏。
可以为候选$v$计算特征$\bm f(h,v)$，把各特征的贡献相加，再归一化：
\begin{equation}
 s_{\bm\theta}(h,v)=\bm\theta^{\mathsf T}\bm f(h,v),\qquad
 p_{\bm\theta}(v\mid h)=\frac{\exp s_{\bm\theta}(h,v)}
 {\sum_{u\in\mathcal V}\exp s_{\bm\theta}(h,u)}.
 \label{eq:nlp-classical-loglinear}
\end{equation}
这就是对数线性参数化：概率的对数在归一化常数之外是特征的线性组合。
同一个特征权重用于所有触发该特征的样本，因此少量参数可以联系许多表格单元。
若全部特征都是独立单元的指示函数，足够自由的有限权重可以表示任意严格正的行概率表，
例如取权重为各项概率的对数。概率为零的表格边界则需要权重趋向负无穷的极限，
或者另加硬掩码排除相应事件；有限softmax本身不会产生精确零概率。
因此，使用特征参数化并不意味着自动获得有用共享。

在式~\eqref{eq:nlp-classical-crf}中，把候选从单个词换成整条标签序列，把特征沿位置相加，
就得到全局归一化的结构评分。局部归一化与全局归一化的分母不同，不能把每一步的softmax
相乘后仍称作同一个CRF。特征也只能使用运行时可用的信息；表达能力强不代表允许读取答案。

\subsection{连续表示与上下文状态共享}
\label{sec:nlp-classical-neural}

人工特征明确规定了哪些事件共享权重，但未规定的相似性仍无法传递。早期神经概率语言
模型把每个词映射为可学习的连续向量，把有限历史的向量拼接后经过非线性映射，再预测
下一词。这使不同但表示接近的历史能够产生相近的输出，词表示与预测参数共同接受似然
训练，而不是先靠人工给词指定语义坐标\scite{nlp-classical-bengio2003}。

为简化记号，令$\bm e(w)$为词嵌入，$\sigma$为逐分量非线性函数。有限窗口网络与递归
网络分别可写成
\begin{align}
 \bm h_t&=\sigma\!\left(\bm W[\bm e(w_{t-m});\ldots;\bm e(w_{t-1})]+\bm b\right),
 \label{eq:nlp-classical-window}\\
 \bm h_t&=\sigma(\bm U\bm h_{t-1}+\bm W\bm e(w_{t-1})+\bm b),\qquad
 p(w_t\mid w_{<t})=\operatorname{softmax}(\bm O\bm h_t+\bm a)_{w_t}.
 \label{eq:nlp-classical-rnn}
\end{align}
两式是对上下文状态的两种选择，不是同时对同一$\bm h_t$施加的等式。窗口网络只读最近
$m$个词；递归网络让历史反复进入同一个状态更新，因而没有显式的固定窗口，但有限维
状态和训练中的梯度传播限制了它保留远处细节的能力。共享向量使平滑成为可学习的函数，
也带来参数之间的耦合和非凸优化问题。

因此，神经模型与统计建模并非对立分类。它仍然指定概率、结构和估计准则，只是把局部表
换成了可共享的连续函数。下一章的Transformer将进一步改变历史信息如何进入计算，
而不是从此取消概率建模的四个决定。

\section{语言模型的估计与求解}
\label{sec:nlp-classical-estimation}

\subsection{似然估计与稀疏事件平滑}
\label{sec:nlp-classical-smoothing}

固定二元结构后，若历史$h$出现$c(h)$次，其中后继$v$出现$c(h,v)$次，最大似然估计为
$c(h,v)/c(h)$。这由带行归一化约束的对数似然最大化得到：拉格朗日一阶条件使每个概率
正比于对应计数，归一化再确定比例。$c(h)=0$时数据没有决定这一行，不能把$0/0$当作零。

在教学语料中，“会议”后的“改至”有两次、“定于”有一次；最大似然概率为$2/3$和$1/3$，
未出现过的后继概率为零。只要测试句用到一个零概率事件，整句似然就变为零。加性平滑
引入$\alpha>0$并给所有$V$个候继保留质量：
\begin{equation}
 \widehat p_\alpha(v\mid h)=\frac{c(h,v)+\alpha}{c(h)+\alpha V}.
 \label{eq:nlp-classical-additive}
\end{equation}
设固定词表为“会议、改至、定于、周五、取消、EOS”六项，取$\alpha=1$，
“会议”之后的“改至”和“定于”概率分别成为$3/9$和$2/9$，其余四项各为$1/9$。这是构造算例，说明概率从已见事件转给未见事件；它不说明加一平滑
就是自然语言中的最优选择。较大词表会分走大量概率，$\alpha$应在独立验证材料上选择。

也可按$p(v\mid h)=\lambda(h)p_{\rm high}(v\mid h)+(1-\lambda(h))p_{\rm low}(v)$
在高阶与低阶分布间插值。两个分布都归一化且$\lambda\in[0,1]$时，结果仍是分布。
低阶分布怎样估计、计数少时怎样回退，会实质改变未见事件的排序；回退不能理解为
简单地给每个零单元填一个相同常数。

\subsection{隐结构边缘化与迭代估计}
\label{sec:nlp-classical-em}

HMM的词序列可见而状态不可见时，完整状态计数不能直接取得。对每个候选前缀定义
前向量$\alpha_t(j)=p(w_{1:t},z_t=j)$，利用链结构得到
\begin{equation}
 \alpha_1(j)=\pi_jB_{j,w_1},\qquad
 \alpha_t(j)=B_{j,w_t}\sum_i\alpha_{t-1}(i)A_{ij}.
 \label{eq:nlp-classical-forward}
\end{equation}
末尾求和得到观测概率。状态数为$K$时，朴素稠密转移的计算量为$O(TK^2)$；
这是把指数条路径中的共享前缀合并计算的结果，并没有把求和改成最大值。
前向--后向计算还能给出各位置状态和相邻状态对的后验概率。

\term{期望最大化}（expectation--maximization, EM）据此交替计算后验期望计数和更新概率表。
在第$r$轮，用旧参数构造
\begin{equation}
 Q(\bm\theta\mid\bm\theta^{(r)})=
 \mathbb E_{z\mid w,\bm\theta^{(r)}}[\log p_{\bm\theta}(w,z)],\qquad
 \bm\theta^{(r+1)}\in\arg\max_{\bm\theta}Q(\bm\theta\mid\bm\theta^{(r)}).
 \label{eq:nlp-classical-em}
\end{equation}
例如状态转移表按后验期望的$i\to j$次数归一化。精确E步与不降低$Q$的M步可以保证观测
对数似然不下降，却不保证找到全局最优，更不保证状态对应预期语义。近似后验、平滑项
和数值误差都会改变需要检查的目标；不能把精确EM的结论原封不动移给任意近似更新。

\subsection{条件目标、先验与泛化检查}
\label{sec:nlp-classical-conditional-learning}

有标注$(x_i,y_i)$时，可以最大化条件对数似然。以线性链CRF为例，一条样本的梯度是
\begin{equation}
 \begin{aligned}
 \nabla_{\bm\theta}\log p_{\bm\theta}(y_i\mid x_i)
 &=\bm F(x_i,y_i)-\mathbb E_{Y\mid x_i,\bm\theta}[\bm F(x_i,Y)],\\
 \bm F(x,y)&=\sum_t\bm f(y_{t-1},y_t,x,t).
 \end{aligned}
 \label{eq:nlp-classical-crf-gradient}
\end{equation}
第一项是参考结构中的特征次数，第二项是模型在所有合法结构下预期的次数。梯度推动
模型减少两者差异；归一化项的导数解释了第二项为何不能省略。线性链结构使后验期望
可用动态规划计算，换成复杂结构则未必仍然可行。

若取零均值、各向同性高斯先验$\bm\theta\sim\mathcal N(\bm0,\sigma^2\bm I)$，
把其对数加入条件对数似然，就相当于扣除$\lVert\bm\theta\rVert_2^2/(2\sigma^2)$。
这种平方范数惩罚表示在有限数据以外偏好较小参数，而不是增加了观测证据。验证集选择惩罚强度，测试集只用于最终比较。
似然、结构准确率和业务代价回答不同问题：局部概率更校准可能不改变最大分标签，
标签准确率更高也可能漏掉一条关键报名条件。第\ref{chap:nlp-evaluation}章将统一评价协议，应用章再给出各输出
对象的具体核验办法。

\section{本章小结}
\label{sec:nlp-classical-summary}

统计语言建模把有限语料转成可用于新输入的规律，依靠的是明确的约束与共享。概率对象
决定要解释什么，依赖结构决定保留什么信息，参数化决定哪些事件互相借用证据，估计准则
再决定如何从观测得到参数。局部链、递归树与交换性混合并不构成通用能力的排名，表格、
特征和连续表示也不是与这些结构平行的任务分类。

我们已经能够解释为什么未见句子仍可能获得非零概率，也能指出有限历史、隐变量解释与
稀疏估计留下的局限。下一章保持概率与学习目标的明确含义，把历史交互改为可学习的
上下文计算，并在更广泛语料上训练共享参数。预训练改变规模与信息组织，仍须回答本章
建立的建模问题。

\paragraph{延展阅读}
连续表示语言模型、条件随机场和潜在狄利克雷分配分别展开参数共享、条件结构预测与
交换性混合。对照阅读时，可以固定随机对象、分解假设和估计准则，比较三类模型共享
证据的方式及其不能表达的依赖，而不把“统计模型”和“神经模型”视为互斥任务类别
\scite{nlp-classical-bengio2003,pgm-energy-models-sutton2012,blei2003}
